\practicechapter{第 10 章实践：x86-64 汇编语法总览}

本章对应第一册第 10 章“x86-64 汇编语法总览”。正文讲寄存器、立即数、内存操作数、AT\&T 与 Intel 风格差异。本章要做的是把反汇编从“看不懂的一堆字符”变成可核对证据：每条指令至少能分清助记符、源操作数、目的操作数、寄存器还是内存。

\practicesection{一、对应原书问题：反汇编是证据，不是装饰}

很多报告会贴一段 \cmd{objdump -d} 输出，然后直接说“可以看到汇编”。这不够。本章要求你选择一个很小函数，例如 \code{add_i32} 或 \code{main} 中的一段，逐行标注：指令地址、机器码字节、助记符、操作数、操作数含义。

这一步要和本卷解析器连接起来。先用 \code{--sections-csv} 找到 \code{.text}，再用反汇编观察 \code{.text} 中的代码。这样汇编不是孤立文本，而是 ELF section 中字节的一种解释。

\practicesection{二、从特例推到规律：语法风格会改变阅读方向}

特例：AT\&T 风格通常写作 \code{movl \%edi, -4(\%rbp)}，源操作数在前，目的操作数在后；Intel 风格常写作 \code{mov DWORD PTR [rbp-4], edi}，目的操作数在前，源操作数在后。两者描述的可能是同一条机器指令，但文本顺序不同。

由此推出规律：读汇编时先确认语法风格。不要看到两个操作数就凭直觉猜方向。对于本章报告，建议固定使用 \cmd{objdump -d -Mintel} 或默认 AT\&T 风格之一，并在报告开头写明。若对照两种风格，必须说明同一条指令如何对应。

\practicesection{三、实践任务：标注五条指令}

执行：

\begin{lstlisting}[language=bash]
objdump -d ./hello-o0 > reports/hello-o0.att.txt
objdump -d -Mintel ./hello-o0 > reports/hello-o0.intel.txt
\end{lstlisting}

从 \code{main} 或 \code{add_i32} 中选择五条连续指令，并按四个字段标注。第一是地址，例如 \code{401126}，它是反汇编视角的代码位置。第二是机器码字节，例如 \code{55}，它是 CPU 实际解码的输入。第三是助记符，例如 \code{push}，它描述操作种类。第四是操作数，例如 \code{rbp}，它说明本条指令访问的是寄存器、立即数还是内存。

如果编译器生成的指令和书中示例不同，不要强行改报告。写明编译器版本、优化等级和构建命令，然后解释当前产物中实际出现的指令。

\practicesection{四、验收：能读出最小语义}

本章交付物是 \filepath{reports/ch10-assembly-syntax.md}。通过标准如下：

\begin{enumerate}
  \item 报告明确使用 AT\&T 还是 Intel 风格。
  \item 至少标注五条指令，每条都分出助记符和操作数。
  \item 至少说明一条寄存器操作和一条内存操作。
  \item 能把反汇编片段和 \code{.text} section 联系起来。
  \item 不根据助记符名称过度推断性能；本章只训练语法和证据定位。
\end{enumerate}

完成本章后，读者应该能把反汇编当成可阅读证据，而不是只把它当作“底层感很强”的截图。
